\chapter{Ajustes del equipo}\label{cap:14-ajustes}

En XFCE todos los ajustes están en \menu{Aplicaciones > Configuración >
Gestor de configuración}. En los demás escritorios: \menu{Sistema >
Preferencias} (MATE), \menu{Preferencias > Configuración de LXQt} (LXQt),
\emph{Configuración} (GNOME) o \emph{Preferencias del sistema} (KDE
Plasma).

\captura{gestor-configuracion}{El Gestor de configuración de XFCE}

\section{Pantalla}

\menu{Configuración > Pantalla} cambia la resolución (el tamaño de todo en
la pantalla) y, si conectas un monitor o un proyector, si se ve lo mismo
en los dos o si se extiende el escritorio.

\begin{consejo}
Si todo se ve demasiado pequeño en una pantalla grande, en lugar de bajar
la resolución sube el tamaño de letra en \menu{Apariencia > Tipos de
letra} y, en \menu{Apariencia > Configuración}, la \emph{escala}.
\end{consejo}

\section{Sonido}

El icono del altavoz de la bandeja sube y baja el volumen. Con clic
derecho, \menu{Configuración del sonido} muestra todas las salidas y
entradas: elige ahí los audífonos, el parlante Bluetooth o el micrófono
para las videollamadas.

\section{Teclado y ratón}

\begin{itemize}
  \item \menu{Configuración > Teclado > Distribución} cambia el idioma del
        teclado. Para usar dos (por ejemplo, español e inglés), agrégalos
        aquí y elige una combinación para cambiar entre ellos.
  \item \menu{Configuración > Ratón y panel táctil}: velocidad del puntero,
        \emph{tocar para hacer clic} y el desplazamiento con dos dedos en el
        touchpad.
\end{itemize}

\section{Fecha y hora}

La hora se ajusta sola por internet. Para cambiar la zona horaria, haz
clic derecho en el reloj del panel, \menu{Propiedades}. Si la hora está mal
después de usar Windows en el mismo equipo, mira la
sección~\ref{sec:prob-hora}.

\section{Idioma}

Lázaro queda en el idioma que elegiste al instalar, y todos los programas
lo siguen. Cambiar el idioma de un sistema ya instalado todavía no tiene un
botón: si lo necesitas, pide ayuda (capítulo~\ref{cap:17-pedir-ayuda}).
Para escribir en otro idioma no hace falta cambiarlo: basta con agregar
otra distribución de teclado (arriba) y, en LibreOffice, elegir el idioma
del texto en \menu{Herramientas > Idioma}.

\section{Usuarios y contraseñas}

\subsection{Cambiar tu contraseña}
Abre una terminal, escribe \texttt{passwd} y pulsa \keys{Enter}. Pide la
contraseña actual y dos veces la nueva. Mientras escribes no se ve nada: es
normal.

\subsection{Entrar sin contraseña al encender}
Lázaro recuerda esta opción en el archivo \texttt{/etc/lazaro/autologin}.
Para dejar de entrar solo, en la terminal:
\begin{terminal}
doas rm /etc/lazaro/autologin
\end{terminal}
Para empezar a entrar solo con tu usuario:
\begin{terminal}
echo "\$USER" | doas tee /etc/lazaro/autologin
\end{terminal}
En los dos casos, reinicia para que tome efecto.

\subsection{Otra persona en el mismo equipo}
Cada persona puede tener su propio usuario, con su carpeta y su
configuración. En la terminal, cambiando \texttt{ana} por el nombre:
\begin{terminal}
doas adduser ana
\end{terminal}

\subsection{Qué es doas}
Algunas órdenes cambian el sistema y necesitan permiso de administrador.
Se escriben con \texttt{doas} delante, que pide tu contraseña. Es lo mismo
que el \emph{sudo} de otras distribuciones o el \enquote{Ejecutar como
administrador} de Windows. Usa \texttt{doas} solo con órdenes que
entiendas o que vengan de una fuente de confianza, como este manual.

\section{Energía y batería}

\menu{Configuración > Gestor de energía} decide cuándo se apaga la
pantalla, qué pasa al cerrar la tapa del portátil y a qué nivel de batería
se suspende. Si el equipo tiene batería, Lázaro instala los \emph{modos de
energía}: en el icono de la batería se elige entre \emph{ahorro},
\emph{equilibrado} y \emph{rendimiento}.

\section{Bluetooth}

En KDE Plasma y GNOME, el Bluetooth está en la configuración del sistema y
en el menú de arriba a la derecha (GNOME) o en la bandeja (Plasma). En
XFCE, MATE y LXQt, instala desde la tienda el \emph{Gestor de Bluetooth}
(Blueman): agrega un icono a la bandeja para buscar y emparejar
audífonos, parlantes, ratones y teclados.

Para emparejar, pon el aparato en modo de emparejamiento (casi siempre,
manteniendo apretado su botón de encendido hasta que parpadee) y elígelo
en la lista.

\fuentes
